\documentclass[11pt,a4paper]{article}
\usepackage[T1]{fontenc}
\usepackage[utf8]{inputenc}
\usepackage{lmodern}
\usepackage{amsmath,amssymb}
\usepackage[provide=*,spanish]{babel}
\usepackage[a4paper,top=1.4cm,bottom=1.5cm,left=1.6cm,right=1.6cm,headheight=13pt,headsep=8pt]{geometry}
\usepackage{booktabs,array,tabularx}
\usepackage[labelsep=colon,font=small,skip=3pt]{caption}
\usepackage{enumitem}
\setlist{nosep,leftmargin=1.2em}
\usepackage{xcolor}
\usepackage{fancyhdr}
\pagestyle{fancy}\fancyhf{}
\fancyfoot[C]{\small\thepage}
\renewcommand{\headrulewidth}{0pt}
\usepackage{titlesec}
\titleformat*{\section}{\large\bfseries}
\titlespacing*{\section}{0pt}{7pt}{3pt}
\setlength{\parskip}{2pt}
\addto\captionsspanish{\renewcommand{\tablename}{Tabla}}

\begin{document}

\begin{center}
{\Large\bfseries Nota para el foro}\\[3pt]
{\normalsize Proyecto final CA-0309 \,|\, Grupo 2, Auditoría o Control \,|\, Ronda 1}\\[2pt]
{\small Andrey Gonzales (C33329), Randal Picado Bermúdez (C36024), Anthonny Flores Rojas (C32975)}
\end{center}

\section*{Escenario}
Prima anual de 55\,000 y renuncia anual de decremento único del 6\,\%, con los demás parámetros de
la base (rendimiento del 5\,\%, gasto de 10\,000, descuento del 4\,\% y el mismo cuadro de reservas
por contrato). Se recalcularon probabilidades competitivas, población, flujos, reservas totales y
recursos. El 6\,\% es un estrés académico, no una estimación de la reacción de los afiliados. Se
compara con nuestro caso base de C; no se combina con el escenario adverso de D y no modifica B ni E.

\section*{Comparación propia}
\begin{minipage}[t]{0.50\textwidth}
\vspace{0pt}
\centering\small
\begin{tabular}{lrr}
\toprule
Indicador & Base & Estrés G2\\
\midrule
Activos al cierre del año 5 & 84\,700.19 & 72\,387.73\\
NPV (millones de u.m.) & $-12\,912.86$ & $-10\,222.77$\\
Margen sobre primas desc. & $-57.88\,\%$ & $-44.09\,\%$\\
Total de aportes (millones) & 3\,050.20 & 0.00\\
$\tilde F_1 - V_1$ (millones) & $-2\,024.61$ & $-1\,307.79$\\
\bottomrule
\end{tabular}
\captionof{table}{Base de C frente al estrés del grupo 2.}
\end{minipage}\hfill
\begin{minipage}[t]{0.47\textwidth}
\vspace{0pt}
\small
\textbf{Efecto conjunto.} La cartera termina con 12\,312 activos menos ($-14.5\,\%$). Las primas
totales suben de 24\,017 a 24\,898 millones y los beneficios bajan de 24\,847 a 22\,983 millones
(373 fallecimientos cubiertos menos). El NPV mejora y desaparecen los aportes, pero el producto
sigue perdiendo por contrato: el siniestro esperado (45\,943 a 56\,512) supera la nueva prima neta
de 42\,250. La mejora viene de cubrir a menos personas, no de un mejor producto. No se atribuyen
importes a cada cambio por separado.
\end{minipage}

\section*{Observaciones de Auditoría o Control}
\small
\begin{tabularx}{\textwidth}{@{}>{\bfseries}p{2.4cm}X@{}}
\toprule
\multicolumn{2}{@{}l}{\normalsize\textbf{1. La tarifa no cubre el siniestro esperado en ningún año}}\\
\midrule
Evidencia & La prima neta de gasto y comisión (37\,500 por contrato) queda por debajo del siniestro
esperado por contrato vigente en los cinco años (46\,659 a 57\,392). NPV base $-12\,912.86$
millones; la adquisición explica 7\,800.30 millones (60\,\%).\\
Consecuencia & La pérdida es estructural y crece con la edad de la cartera (edad media de 51.2
años); persiste en el escenario adverso ($-13\,617.35$) y en nuestro estrés ($-10\,222.77$).\\
Acción & Recalcular la prima con la mortalidad de esta población y revisar el gasto de adquisición
antes de aprobar; exigir NPV no negativo en base y adverso.\\
\midrule
\multicolumn{2}{@{}l}{\normalsize\textbf{2. Los recursos no respaldan la reserva desde el primer cierre}}\\
\midrule
Evidencia & $\tilde F_1 - V_1$ es negativo en los tres escenarios: $-2\,024.61$ (base),
$-2\,197.78$ (adverso) y $-1\,307.79$ (estrés). Los recursos iniciales son 25\,000 por contrato y la
reserva del año 1, 40\,000. La base requiere 3\,050.20 millones de aportes.\\
Consecuencia & Aunque el estrés elimina los aportes, su respaldo sigue negativo del año 1 al 4: la
falta de respaldo no se corrige con más renuncias.\\
Acción & Condicionar cualquier aprobación a un plan de capital aprobado que asegure
$\tilde F_t \ge V_t$ en cada cierre, empezando por el año 1.\\
\bottomrule
\end{tabularx}
\normalsize

\section*{Preguntas para el foro}
\begin{enumerate}[label=\arabic*.]
\item \textbf{Al grupo 1 (Actuarial):} con la prima de 50\,000, después de gastos y comisión quedan
37\,500 por contrato, pero el siniestro esperado va de 46\,659 a 57\,392 en los cinco años. Según su
análisis, ¿qué prima haría falta para que el producto no pierda dinero y cómo la justifican?
\item \textbf{Al grupo 7 (Comité de Riesgos):} ¿cuánto capital estarían dispuestos a poner para
cubrir la reserva desde el primer año (al menos 2\,025 millones en la base) y qué cifra mirarían
para decidir si el producto sigue?
\end{enumerate}

\section*{Opinión preliminar}
Posponer el seguro colectivo hasta contar con una tarifa que dé NPV no negativo en base y
adverso y con un plan de capital que respalde la reserva desde el año 1. Apoyar la cuenta
equity-linked con $G = 100\,000$ (NPV medio 4\,054.40, frecuencia de pérdida 24.6\,\%, CTE95
22\,959.79), siempre que se vigilen las pérdidas más grandes.

\noindent\textit{Límites:} son cifras del ejercicio, no de solvencia regulatoria; el 3\,\% de renuncia no viene
de datos observados y no calculamos FCF, RA ni CSM. La opinión final se da en el cierre oral.

\end{document}